\documentclass[10pt,a4paper]{amsart}
\usepackage[margin=19mm]{geometry}
\usepackage{amsmath,amssymb,mathtools}
\usepackage[scr=rsfs,cal=euler]{mathalfa}
\usepackage{xfrac}
\usepackage[unicode,hidelinks,bookmarks=false]{hyperref}
\newcommand{\ie}{i.e.\ }
\newcommand{\cf}{cf.\ }
\newcommand{\resp}{respectively\ }
\newcommand{\Subsection}{\subsection}
\providecommand{\nobreakdash}{\nobreak\hbox{-}}
\setlength{\emergencystretch}{2em}
\multlinegap0pt
\usepackage{braket}
\usepackage{amssymb}
\usepackage{float}
\usepackage{xcolor}
\usepackage{algorithm}\usepackage{algpseudocode}
\usepackage{tikz}
\usepackage{bm}
\usepackage{mathtools}
\usepackage{dsfont}
\usepackage{csquotes}
\newtheorem{lemma}{Lemma}
\title{The Power Contamination Problem on Grids Revisited: Optimality, Combinatorics, and Links to Integer Sequences}
\author{El-Mehdi Mehiri$^{a,b}$ and Mohammed L. Nadji$^{b}$}
\date{Source: arXiv:2509.12756v3 (CC BY 4.0)}
\begin{document}
\maketitle

\noindent\emph{Source and licence.} The Power Contamination Problem on Grids Revisited: Optimality, Combinatorics, and Links to Integer Sequences. Version arXiv:2509.12756v3 (CC BY 4.0), distributed by arXiv under the Creative Commons Attribution 4.0 International licence, http://creativecommons.org/licenses/by/4.0/, as recorded in the arXiv licence metadata for this exact version and shown on the abstract page https://arxiv.org/abs/2509.12756v3. Full licence notice: \texttt{licenses/arxiv-2509.12756-CC-BY-4.0.txt}.\par\smallskip

\noindent\textit{Editorial scope.} Counted unit: the article's lemma lem:3row-avoid (source0.tex lines 3385--3394). The article's complete original proof of it is delivered with it (source0.tex lines 3396--3449, one \texttt{proof} environment of 54 lines). No mathematics is written, repaired or re-derived: the statement and the proof are the article's own lines, in the article's own notation, and no retained line is substituted or re-flowed.  No further source-local context is needed: every cross-reference of every retained line resolves inside the two delivered spans.  Nothing was omitted from any retained span, so the package carries no omission record.  No retained line contains a citation, so the wrapper carries no bibliography and no bibliography entry.  No retained line contains a figure environment, an image-inclusion command or drawing code, so no image is delivered and the package declares no asset dependency.  Editorial formatting only, in addition: an AMS \texttt{amsart} preamble that restates the article's own declarations and macros used by the retained lines, namely the environments \texttt{lemma} on separate counters, together with the standard packages those lines need.  The retained environments print in this document as Lemma 1 in order of appearance, whereas the article prints the same items with the numbering of its own full document.  The complete native source of the article as distributed in the arXiv e-print for this exact version is bundled unchanged as \texttt{sources/185134/source0.tex}.
\begin{lemma}\label{lem:3row-avoid}
Let \(S\) be an optimal contamination set of \(G(3,2k+1)\), and let
\(w=r_0r_1\cdots r_k\in\{1,2,3\}^{k+1}\) be its associated word. If \(w\)
avoids both factors \(13\) and \(31\), then
\[
C(S)=D(S).
\]
In particular, each column of \(C(S)\) contains at most one contaminated cell,
and therefore \(S\) does not fully contaminate \(G(3,2k+1)\).
\end{lemma}

\begin{proof}
We first prove that \(D(S)\subseteq C(S)\). Every cell of \(S\) is contaminated
initially. Now let \(1\le j\le k\) and assume that \(r_{j-1}=r_j\). Then the
two contaminated cells $(r_{j-1},2j-1)$ and $(r_j,2j+1)$ lie in the same row, so by rule~\textup{(\texttt{d})} the cell \((r_j,2j)\) becomes
contaminated. Hence every cell of \(D(S)\) belongs to \(C(S)\).

It remains to show that no cell outside \(D(S)\) can be contaminated from
\(D(S)\). For this, we record three immediate observations.

 
\noindent\textit{(i) Each column contains at most one cell of \(D(S)\).}
Indeed, each odd column contains exactly one cell of \(S\), and an even column
contains at most one cell, namely \((r_j,2j)\) when \(r_{j-1}=r_j\).

 
\noindent\textit{(ii) If two consecutive columns are both nonempty in \(D(S)\),
then their contaminated cells lie in the same row.}
Indeed, if column \(2j\) is nonempty, then by definition its unique cell is
\((r_j,2j)\), while column \(2j+1\) contains \((r_j,2j+1)\). The same argument
applies to the pair \((2j-1,2j)\).

 
\noindent\textit{(iii) If two columns at distance \(2\) are both nonempty in
\(D(S)\), then their contaminated cells lie either in the same row or in two
consecutive rows.}
If the two columns are odd, say \(2j-1\) and \(2j+1\), then their contaminated
cells are \((r_{j-1},2j-1)\) and \((r_j,2j+1)\). Since \(w\) avoids \(13\) and
\(31\), we have \(|r_j-r_{j-1}|\le 1\). If the two columns are even, say
\(2j\) and \(2j+2\), then both being nonempty implies $r_{j-1}=r_j=r_{j+1}$,  so the two contaminated cells lie in the same row.

 
We now check that none of the rules \textup{(\texttt{a})}--\textup{(\texttt{h})} can produce a
cell outside \(D(S)\).

Rules \textup{(\texttt{a})} and \textup{(\texttt{b})} require two contaminated cells in columns
at distance \(2\) and in rows \(1\) and \(3\), which is impossible by
observation~(iii). Rule \textup{(\texttt{c})} requires two contaminated cells in the
same column, which is impossible by observation~(i). Rules
\textup{(\texttt{e})}--\textup{(\texttt{h})} require two contaminated cells in consecutive
columns and adjacent rows, which is impossible by observation~(ii).

Finally, consider rule~\textup{(\texttt{d})}. If it applies to two contaminated cells in
the same row and in columns at distance \(2\), then either the target cell lies
in an even column \(2j\), in which case it is exactly \((r_j,2j)\in D(S)\), or
the target cell lies in an odd column, which already contains its unique
initial contaminated cell from \(S\subseteq D(S)\).

Therefore no cell outside \(D(S)\) can be contaminated from \(D(S)\). Since
\(D(S)\subseteq C(S)\), it follows that \(C(S)=D(S)\).

By observation~(i), each column of \(C(S)\) contains at most one contaminated
cell. Hence \(C(S)\neq V(G(3,2k+1))\), so \(S\) does not fully contaminate the
grid.\qedhere
\end{proof}

\end{document}
